\labsessionheader{2026-09-20}{Feature branches archived as tags}%
  {5958d9e5}{n/a}{bd gt1z; archive manifest in figures/021/}

\subsection*{Result}
After the owner requested further cleanup, \bd{gt1z} archived the remaining
39 local feature branch tips as tags under
\texttt{archive/2026-09-20/}. The tags are available locally and on origin.
All 39 local feature branches and their 22 remote counterparts were removed.
The remaining branches are \texttt{main} locally, and \texttt{main} plus
\texttt{\_\_dolt\_remote\_info\_\_} on origin. The metadata branch supports
Beads synchronization.

\subsection*{Verification}
\texttt{git fetch --prune origin} refreshed the inventory, and
\texttt{gh pr list --state open --limit 100} again returned no open pull
requests. All feature branch tips matched entry 020's inventory;
\texttt{git worktree list --porcelain} confirmed none was checked out.
\texttt{git bundle verify} rechecked the earlier recovery bundle, and
\texttt{git bundle list-heads} confirmed that it contained every candidate tip.

One atomic push created the 39 archive tags and removed the 22 remote
branches. Each deletion required the expected old tip, and each tag required
that its name did not already exist. After the push,
\texttt{git ls-remote origin} verified every archive tag's exact commit ID
and the absence of the deleted remote branches. Only then were the local
branch names removed. The mapping is recorded in
\texttt{figures/021/archive-manifest.json}.

No compiler tests were run because no source, fixture, golden, or harness
changed. Existing detached worktrees and uncommitted files were retained.
The lab-report manifest and page budget were checked before sealing, as
required by \texttt{CLAUDE.md}'s session-close procedure.

\subsection*{Recovery}
For example, restore the former \texttt{3qlq-inc1} branch with:
\begin{verbatim}
git switch -c 3qlq-inc1 archive/2026-09-20/3qlq-inc1
\end{verbatim}
The original bundle remains at
\texttt{.git/branch-cleanup-20260920/before.bundle}.
Archiving preserves the divergent histories; it does not establish that
their changes were merged into main.

\subsection*{Update to entry 020}
Entry 020 recorded these branches as retained pending review. This follow-up
resolves their storage and branch-list cleanup through remote archive tags.
